\documentclass[14pt,aspectratio=169]{beamer}
\usetheme{Copenhagen}
\usecolortheme{beaver}
\setbeamertemplate{navigation symbols}{}
% \setbeamercovered{transparent}

\usepackage[utf8]{inputenc}

\title{Type-checking Linearity in Haskell's\\ Core Intermediate Language}
\author{Rodrigo Mesquita and Bernardo Toninho}
\institute{\small NOVA School of Science and Technology}
\date{\small February 2023}

\begin{document}
\begin{frame}
    \titlepage
    \begin{figure}[htpb]
        \begin{center}
            % \vspace*{-0.5cm}
            \includegraphics[width=0.4\linewidth]{../logo_nova.png}
        \end{center}
    \end{figure}
\end{frame}

\section{Introduction}

\begin{frame}{Context}

\begin{itemize}
    \item<1-> Statically typed programs eliminate errors at compile time
    \item<2-> Linear types can enforce resource usage invariants
        % Are expressive types
    \item<3-> Haskell is a lazy, purely functional programming language
    \item<4-> Linear types were introduced in Haskell, 31 \emph<4>{years} after Haskell's inception.
                % recently introduced
                % Which came with new challenges such as backwards compatibility, future
                % proofing and code re-use
\end{itemize}

\end{frame}

\begin{frame}{Motivation}
\begin{itemize}
    \item<1-> However, linear types were also introduced in \textbf<1>{Core}
        % a typed language that serves as Haskell's main intermediate representation
        % Core is explicitly typed
        % Core is very important because it serves as an internal consistency
        % tool and validation layer of the implementation
    \item<2-> Core programs undergo a myriad of optimizing transformations
        % These (core-to-core) transformations follow a modular approach, where
        % each transformation produces Core that is further used as input to
        % other transformations
        % Which might introduce lazy let bindings (through e.g. beta reduction),
        % inline definitions, move bindings around (let-floating
        % transformations), 
    \item<3-> \textbf<3>{Core's linear type-checker doesn't accept linear
        programs} after optimizing transformations are applied
        % Optimizing transformations destroy linearity as the type-checker sees
        % it. However, we believe optimizing transformations *don't* destroy
        % linearity, it's just a limitation of the type system
    \item<4-> \textbf<4>{Core's linear type-checker is disabled} because
        otherwise disabling optimizations is far worse
        % it doesn't accept transformed Core programs
        % GHC's Core type-checker is disabled because it doesn't accept Core
        % programs that undergo optimizing transformations in Core
        % and otherwise disabling optimizing transformations is not feasible.
\end{itemize}
\end{frame}

\begin{frame}{Objectives}
\begin{itemize}
    \item<1-> Extend Core's type-system to accommodate more linear Core programs
        % Big distinction between syntatically valid linearity and semantically
        % valid. Transformations result in many semantically valid but
        % syntatically invalid programs, and we need to extend Core to
        % account for this semantic linearity.
    \item<2-> Validate optimizing transformations with the improved Core linear type-system
    \item<3-> Implement the extensions to Core in GHC, the leading Haskell
        compiler
\end{itemize}
\end{frame}

\section{Background}

\begin{frame}{}
\end{frame}

\end{document}
